% ECONOMICS_PROBLEM: {"id":"C73-15","title":"Payoff characterization of convergence for a fixed mirror-descent rule","jel_code":"C73","source_id":"OP-0148"}
% FIRST_POSTED: 2026-10-10
% ATTEMPT_STATUS: solved
% ENTRY_KIND: research_attempt
\documentclass[11pt,a4paper]{article}
\usepackage[left=20mm,right=20mm,top=18mm,bottom=18mm]{geometry}
\usepackage[T1]{fontenc}
\usepackage{lmodern,microtype}
\usepackage{amsmath,amssymb,amsthm,mathtools}
\usepackage{xurl,hyperref}
\hypersetup{hidelinks,hypertexnames=false,pdftitle={C73-15: A closed-certificate characterization for every fixed mirror-descent rule}}
\setlength{\parindent}{0pt}
\setlength{\parskip}{3pt}
\setlength{\emergencystretch}{2em}
\allowdisplaybreaks
\raggedbottom
\widowpenalty=10000
\clubpenalty=10000
\displaywidowpenalty=10000
\newtheorem{theorem}{Theorem}
\newtheorem{lemma}[theorem]{Lemma}
\newtheorem{proposition}[theorem]{Proposition}
\newcommand{\R}{\mathbb R}
\newcommand{\N}{\mathbb N}
\newcommand{\NE}{\operatorname{NE}}
\newcommand{\ri}{\operatorname{relint}}
\title{C73-15: A closed-certificate characterization\\for every fixed mirror-descent rule}
\author{Ji Ho Bae}
\date{10 October 2026}
\begin{document}
\maketitle

\section*{Model, process, and claim status}
\textbf{Model:} GPT Codex. \textbf{Contributor:} Ji Ho Bae.
The exact serving revision was not exposed. The mathematical argument was
produced with parallel derivation and adversarial checking in response to a
request for a complete hand proof of the exact catalogue statement.
No human-written mathematical edits are included. This is an LLM claim of a
complete necessary-and-sufficient characterization for the original fixed-rule
target; community expert review is pending. No proof-assistant verification
is claimed.
The complete original generated output is preserved at the immutable
\href{https://github.com/jbaelaw/fixed-mirror-convergence-characterization/blob/000960f40cf5f7565032a9ba5a2986c93b5e0c5b/paper/Proof.tex}{manuscript TeX}
and \href{https://github.com/jbaelaw/fixed-mirror-convergence-characterization/blob/000960f40cf5f7565032a9ba5a2986c93b5e0c5b/paper/Proof.pdf}{manuscript PDF}.

The characterization is infinitary. No finite recognition algorithm,
step-size-independent potential criterion, or familiar-game-class list is
claimed. The construction retains the arbitrary specified regularizers and
positive schedule, including allowed boundary iterates. The fixed
mirror-descent family, schedule, and universal-interior-initialization target
are the editorial specialization C73-15 of the broader source agenda; the
result does not classify every learning dynamic in that agenda.

\section{Exact target and the form of the answer}

Fix \(d\ge1\) players and nonempty finite action sets \(S_i\), \(i=1,\ldots,d\),
\(X_i=\Delta(S_i)\), and \(X=\prod_iX_i\).
The fixed rule has proper lower-semicontinuous strictly convex regularizers
\(h_i:X_i\to\R\cup\{+\infty\}\), finite and differentiable on
\(\ri X_i\), and steps \(\eta_t>0\). For a game's payoff array \(U\),
write \(u_i^U\) for its multilinear extension and
\(v_i^U(x)_a=u_i^U(a,x_{-i})\). The exact simultaneous update is
\begin{equation}
x_i^{t+1}=\underset{p\in X_i}{\arg\max}
\left\{\eta_t\langle v_i^U(x^t),p\rangle
 -h_i(p)+h_i(x_i^t)
 +\langle\nabla h_i(x_i^t),p-x_i^t\rangle\right\}.
\label{eq:md}
\end{equation}
We retain the canonical assumption that these updates are uniquely defined and
remain in the declared gradient domain from every interior initialization.
No assertion is made about a payoff array for which this assumed rule is undefined.

Our answer associates to \(U\) countable trees \(T_{a,b}(U)\),
\(a,b\in\N_{>0}\). Node validity is a finite closed feasibility problem:
it uses fixed regularizer epigraphs, regularizer evaluations at interior probes
and interpolated interior points, the specified steps, and explicit payoff
polynomials. The interpolated points depend on the finite feasibility variables;
the regularizers themselves remain fixed rule data. It contains no Nash set,
convergence assertion, infinite trajectory, or gradient evaluation.
The condition on the payoff array is
\begin{equation}
\boxed{\quad T_{a,b}(U)\text{ is well-founded for every }a,b\ge1.\quad}
\label{eq:wf}
\end{equation}
Equivalently each tree admits a strictly decreasing countable-ordinal ranking.
The full definition is in Section~4. Sections~2--3 explain why its finite
constraints encode interior-ray slopes and exact updates without imposing
Fr\'echet differentiability or continuity of the boundary gradient.

\section{Closed certificates for interior-ray slopes}

For each \(i\), fix an enumeration \((q_{i,\ell})_{\ell\ge1}\) of the
rational points of \(\ri X_i\). It is dense in \(\ri X_i\), hence in \(X_i\).
All numbers \(h_i(q_{i,\ell})\) are finite and fixed before testing \(U\).
For a singleton simplex repeat its unique point.

We normalize a declared gradient by choosing its representative \(g\) with
\(\sum_{s\in S_i}g_s=0\). Adding a constant to all gradient coordinates
has no effect on a Bregman divergence or on \eqref{eq:md}.
The proof uses only the defining directional property of a gradient on each
interior ray:
\begin{equation}
\lim_{\theta\downarrow0}
\frac{h_i(x+\theta(q-x))-h_i(x)}{\theta}
=\langle\nabla h_i(x),q-x\rangle\qquad(q\in\ri X_i).
\label{eq:dp}
\end{equation}
This holds both for a Fr\'echet gradient and for a linear finite-domain
G\^ateaux gradient. We do not assert a uniform expansion at a boundary point,
or finite regularizer values in a whole boundary neighborhood. A selected
subgradient unrelated to these directional derivatives would be a different
notion from the declared gradient.

\begin{lemma}[Dense interior-ray certificate]\label{rays}
Let \(P\) be a finite-dimensional simplex and let
\(h:P\to\R\cup\{+\infty\}\) be proper, lower semicontinuous and convex,
finite on \(\ri P\). Fix dense interior probes \((q_\ell)\).
For \(x\in P\), \(c\in\R\), and a normalized vector \(g\),
the following are equivalent:
\begin{enumerate}
\item \(h(x)=c\), and the derivative along each fixed probe ray exists and satisfies
\begin{equation}
\lim_{\theta\downarrow0}\frac{h(x+\theta(q_\ell-x))-h(x)}{\theta}
=\langle g,q_\ell-x\rangle\qquad(\ell\ge1).
\label{eq:d}
\end{equation}
\item \(c\ge h(x)\), and, for every \(\ell\),
\begin{equation}
E_\ell:=h(q_\ell)-c-\langle g,q_\ell-x\rangle\ge0;
\label{eq:g1}
\end{equation}
for every pair of integers \(\ell,k\ge1\) there is an integer
\(N_{\ell,k}\ge1\) such that, writing
\(z_{\ell,k}=(1-1/N_{\ell,k})x+q_\ell/N_{\ell,k}\),
\begin{equation}
0\le h(z_{\ell,k})-c-
\frac{\langle g,q_\ell-x\rangle}{N_{\ell,k}}
\le\frac{\|q_\ell-x\|}{kN_{\ell,k}}.
\label{eq:g2}
\end{equation}
\end{enumerate}
The conclusion is the displayed ray identity. It does not assert full
Fr\'echet differentiability at an arbitrary coded boundary point.
\end{lemma}

\begin{proof}
A convex function finite on a relative open set is continuous there.
One elementary justification is to enclose a small neighborhood in a simplex
whose vertices lie in that set. Convexity bounds the function above by the
largest vertex value; reflection about the center then bounds it below.
Applying these bounds on a slightly larger neighborhood gives bounded secant
slopes, hence continuity on the smaller one.

Assume the second condition. The epigraph inequality makes \(h(x)\) finite.
Interior continuity extends \eqref{eq:g1} to all interior \(q\). Choose one fixed
interior \(q^0\) and put \(z_\theta=(1-\theta)x+\theta q^0\), \(0<\theta\le1\).
Convexity and lower semicontinuity give
\[
\lim_{\theta\downarrow0}h(z_\theta)=h(x):
\quad h(z_\theta)\le(1-\theta)h(x)+\theta h(q^0),
\quad h(x)\le\liminf_{\theta\downarrow0}h(z_\theta).
\]
Letting \(\theta\downarrow0\) in \eqref{eq:g1} at \(z_\theta\) yields \(c\le h(x)\).
Thus \(c=h(x)\).

The supporting inequality extends to any finite-value boundary point \(p\)
by applying it to \((1-\theta)p+\theta q^0\) and using the same limit identity.
It is automatic when \(h(p)=+\infty\). Consequently
\begin{equation}
h(p)-h(x)-\langle g,p-x\rangle\ge0\qquad(p\in P).
\label{eq:sup}
\end{equation}
% Browser-only layout normalization: omit the samepage wrapper.
% All enclosed mathematical content and the immutable original manuscript are preserved.
For a fixed probe \(q_\ell\), convexity makes its secant quotient
\[
\phi_\ell(\theta)=\frac{h(x+\theta(q_\ell-x))-h(x)}{\theta},
\qquad 0<\theta\le1,
\]
nondecreasing as a function of \(\theta\). Its limit at zero is therefore its
infimum, possibly an extended real number before the certificate is used.
By \eqref{eq:sup} every quotient is at least \(\langle g,q_\ell-x\rangle\).
Multiplying \eqref{eq:g2} by \(N_{\ell,k}\) gives
\[
\phi_\ell(1/N_{\ell,k})
\le\langle g,q_\ell-x\rangle+\|q_\ell-x\|/k.
\]
Hence the infimum is at most the right side for every \(k\), and equals
\(\langle g,q_\ell-x\rangle\). This proves \eqref{eq:d}, without uniformity
over directions.

Conversely, \eqref{eq:d} and secant monotonicity imply \eqref{eq:g1}. For each fixed
\(\ell,k\), the quotients along \(\theta=1/N\) decrease to the finite value
in \eqref{eq:d}, so some integer \(N\) makes their excess at most
\(\|q_\ell-x\|/k\). When \(q_\ell=x\) the excess is identically zero.
This is \eqref{eq:g2}.
\end{proof}

At a point where the original rule declares a gradient, \eqref{eq:dp} and \eqref{eq:d} identify
the certificate vector with its normalized representative. Indeed their
inner products with every \(q_\ell-x\) agree. By density these vectors span
the affine tangent space, and normalization removes the only ambiguity,
an all-coordinate constant. We will use this identification only at points
already identified as actual iterates of the declared rule.

\begin{lemma}[Dense-probe optimizer certificate]\label{opt}
Suppose \(p\in P\), \(c=h(p)\in\R\), and \(y\) is a finite vector.
Then \(p\) maximizes \(\langle y,z\rangle-h(z)\) over \(P\)
if and only if
\begin{equation}
h(q_\ell)-c-\langle y,q_\ell-p\rangle\ge0\qquad(\ell\ge1).
\label{eq:o}
\end{equation}
\end{lemma}
\begin{proof}
The inequalities extend first to the interior by continuity, then to finite-value
boundary points by the convex-segment limit used above. Infinite-value points
cannot improve the objective. The extended inequalities are exactly the
maximization assertion. The reverse implication is immediate.
\end{proof}

These lemmas do not assume continuity of the boundary gradient. For a finite
list of probes, \eqref{eq:g1} is closed. With an integer \(N_{\ell,k}\) fixed,
the interpolated point lies in the relative interior for every \(x\in P\).
Every coordinate of that point is at least the corresponding positive
coordinate of \(q_\ell/N_{\ell,k}\), so all such points belong to a compact
subset of the relative interior. Thus \(h(z_{\ell,k})\) is continuous
in \(x\), and both inequalities in \eqref{eq:g2} are closed.
The epigraph condition \(c\ge h(x)\) is closed by lower semicontinuity.
These closure facts are the reason for retaining ray-specific denominators as
integer node labels instead of assuming that finite mirror-descent trajectories
depend continuously on their initialization.

\section{Payoff polynomials and compact initial cores}

Define pure-deviation regrets
\begin{equation}
R_{i,s}^U(x)=u_i^U(s,x_{-i})-u_i^U(x),\qquad
R^U(x)=\max_{i,s}R_{i,s}^U(x).
\label{eq:r}
\end{equation}
They are explicit polynomials in mixed-action probabilities, with coefficients
from \(U\). For example
\begin{equation}
u_i^U(x)=\sum_{z\in\prod_jS_j}U_i(z)\prod_jx_{j,z_j},
\qquad
v_i^U(x)_s=\sum_{z_{-i}}U_i(s,z_{-i})\prod_{j\ne i}x_{j,z_j}.
\label{eq:p}
\end{equation}
For an integer \(a\ge1\), let
\begin{equation}
K_a=\{x\in X:x_{i,s}\ge1/a\text{ for every }i,s\}.
\label{eq:k}
\end{equation}
Each \(K_a\) is compact, possibly empty, and
\(\bigcup_{a\ge1}K_a=\ri X\). Thus compact cores cover precisely all
permitted initial profiles, even arbitrarily close to the boundary.

\section{The finite node test and the payoff criterion}

Fix \(a,b\ge1\). A node of depth \(r\ge1\) has only finite discrete labels:
\begin{itemize}
\item strictly increasing integer times \(1\le T_1<\cdots<T_r\);
\item player/action labels \((i_j,s_j)\), one for each marked time;
\item positive integer bounds \(B_t\), \(1\le t\le T_r\);
\item positive integer denominators \(N_{i,t,\ell,k}\),
\(1\le i\le d\), \(1\le t\le T_r\), \(1\le\ell,k\le r\).
\end{itemize}
The empty node is the root. Truncation from depth \(r\) to depth \(r-1\)
retains the earlier marked times and their labels, the bounds with
\(t\le T_{r-1}\), and denominator labels with \(t\le T_{r-1}\), \(\ell,k\le r-1\).
Extensions must retain all those labels unchanged. Therefore the set of
possible nodes is countable. The continuous variables below are feasibility
witnesses, not node labels.

A depth-\(r\) node is valid if the following finite system has a solution
in variables \(x^t\in X\), \(c_{i,t}\in\R\),
\(g_{i,t}\in\R^{S_i}\), \(1\le t\le T_r\):
\begin{enumerate}
\item \textbf{Compact bounds and initial core:}
\begin{equation}
x^1\in K_a,\qquad |c_{i,t}|\le B_t,\qquad
|g_{i,t,s}|\le B_t,\qquad \sum_sg_{i,t,s}=0.
\label{eq:f1}
\end{equation}
\item \textbf{Epigraph and finite supporting probes:} for every \(i,t\),
\begin{equation}
c_{i,t}\ge h_i(x_i^t),\qquad
E_{i,t,\ell}:=h_i(q_{i,\ell})-c_{i,t}
-\langle g_{i,t},q_{i,\ell}-x_i^t\rangle\ge0
\quad(1\le\ell\le r).
\label{eq:f2}
\end{equation}
\item \textbf{Finite interior-ray secants:} for every \(i,t\) and
\(1\le\ell,k\le r\), put
\[
z_{i,t,\ell,k}=
\left(1-\frac1{N_{i,t,\ell,k}}\right)x_i^t
+\frac{q_{i,\ell}}{N_{i,t,\ell,k}}.
\]
Require
\begin{equation}
0\le h_i(z_{i,t,\ell,k})-c_{i,t}
-\frac{\langle g_{i,t},q_{i,\ell}-x_i^t\rangle}{N_{i,t,\ell,k}}
\le\frac{\|q_{i,\ell}-x_i^t\|}{kN_{i,t,\ell,k}}.
\label{eq:f3}
\end{equation}
\item \textbf{Finite update probes:} for \(1\le t<T_r\), every \(i\),
and \(1\le\ell\le r\),
\begin{equation}
h_i(q_{i,\ell})-c_{i,t+1}
-\left\langle g_{i,t}+\eta_t v_i^U(x^t),
q_{i,\ell}-x_i^{t+1}\right\rangle\ge0.
\label{eq:f4}
\end{equation}
\item \textbf{Marked payoff gaps:}
\begin{equation}
R_{i_j,s_j}^U(x^{T_j})\ge1/b\qquad(1\le j\le r).
\label{eq:f5}
\end{equation}
\end{enumerate}
The tree \(T_{a,b}(U)\) consists of its root and its valid nodes under the
specified truncation. Validity is inherited by truncation, since all retained
constraints occur among the child's constraints.

Every test is finite and closed in its real feasibility variables. In particular,
for each fixed denominator label the argument of the regularizer in \eqref{eq:f3}
ranges in a compact subset of its relative interior, where it is continuous. The payoff dependence occurs only in
the polynomials \eqref{eq:p}, \eqref{eq:f4}, and \eqref{eq:f5}. Everything else belongs to the fixed rule
and the fixed probe system. The root can be included even when \(K_a\) is empty;
then it has no valid children and is well-founded.

\begin{theorem}[Full payoff characterization]\label{main}
For every payoff array \(U\) for which the canonical fixed rule is well-defined
and gradient-domain preserving as assumed, the following are equivalent:
\begin{enumerate}
\item For every \(x^1\in\ri X\), the exact simultaneous rule \eqref{eq:md} satisfies
\(\operatorname{dist}(x^t,\NE(U))\to0\).
\item The payoff-array criterion \eqref{eq:wf} holds for the finite tests \eqref{eq:f1}--\eqref{eq:f5}.
\item For every \(a,b\), there is a countable ordinal \(\alpha_{a,b}\) and
a map \(\rho_{a,b}:T_{a,b}(U)\to\alpha_{a,b}\) strictly decreasing
from a node to each proper extension.
\end{enumerate}
\end{theorem}

\section{Proof of the full equivalence}

\textbf{Regret and distance have the same vanishing criterion.}
The function \(R^U\) is continuous and nonnegative: a player's mixed payoff
is a convex combination of its pure-action payoffs, so at least one of its
pure deviations has nonnegative regret. Moreover \(R^U(x)=0\) if and only
if every pure deviation is unprofitable, which is equivalent by linearity
to \(x\in\NE(U)\). The Nash set is nonempty by the finite-game existence
theorem \cite{nash} and is closed in compact \(X\).

For any sequence \(z_t\in X\), if its distance to the Nash set tends to zero,
uniform continuity of \(R^U\) implies \(R^U(z_t)\to0\). Conversely, if
\(R^U(z_t)\to0\) but its Nash distance does not, a subsequence with distance
at least some \(\varepsilon>0\) has a compact subsequential limit \(z\).
Continuity gives \(R^U(z)=0\), hence \(z\in\NE(U)\), a contradiction.

\textbf{Failure of convergence produces an infinite branch.}
Suppose an actual trajectory from interior \(x^1\) fails to approach the Nash
set. Choose \(a\) with \(x^1\in K_a\). Its nonnegative regrets fail to
tend to zero, so for some \(b\ge1\) there are strictly increasing times
\(T_j\) with \(R^U(x^{T_j})\ge1/b\). Choose a player and a pure action
attaining this maximum at each such time.

Put \(c_{i,t}=h_i(x_i^t)\), and let \(g_{i,t}\) be the normalized relative
gradient supplied by the declared rule. All these numbers are finite.
Choose integer \(B_t\) exceeding their absolute values. Property \eqref{eq:dp} and
Lemma~\ref{rays} supply integer denominators \(N_{i,t,\ell,k}\) for every date,
probe, and precision.
Fix each such choice once. Lemma~\ref{opt} applied to \eqref{eq:md} gives all the
update inequalities. At depth \(r\), the actual finite prefix through \(T_r\)
satisfies \eqref{eq:f1}--\eqref{eq:f5}. Its labels extend the preceding node. These nodes form
an infinite branch of \(T_{a,b}(U)\), so \eqref{eq:wf} fails.

\textbf{An infinite branch produces one exact bad trajectory.}
Conversely, suppose \(T_{a,b}(U)\) has an infinite branch. Its marked times
strictly increase and therefore tend to infinity. Along the branch every
bound \(B_t\) is eventually assigned and never changed. Each denominator
\(N_{i,t,\ell,k}\) is also eventually assigned and never changed.

At each depth choose one finite feasibility witness. Extend it beyond its last
date by arbitrary points of the coordinate boxes
\begin{equation}
\prod_{t\ge1}\left[X\times
\prod_i[-B_t,B_t]\times\prod_i[-B_t,B_t]^{S_i}\right].
\label{eq:box}
\end{equation}
These boxes are all compact and nonempty. Extract a subsequence converging in
the first date's coordinates, then a subsequence converging in the second date's
coordinates, and so on. The diagonal subsequence converges coordinatewise to
arrays \((x^t,c_{i,t},g_{i,t})_{t\ge1}\). Each constraint from any fixed node
involves only finitely many coordinates and is closed; all sufficiently deep
witnesses satisfy it. Its limit therefore satisfies that constraint. This proves
simultaneously every constraint on the branch, without assuming continuity
of a mirror-descent trajectory map.

For fixed \(i,t\), each probe \(\ell\) and precision \(k\) eventually occurs
in a branch node. Hence the limit arrays satisfy the full conditions of
Lemma~\ref{rays}, with the same fixed denominator for each pair. It follows
that \(c_{i,t}=h_i(x_i^t)\), and the vector \(g_{i,t}\) has the certified
linear slopes on every fixed interior probe ray. We do not infer a full
Fr\'echet gradient at an arbitrary coded boundary point.
All update probes eventually occur as well, so Lemma~\ref{opt} identifies
\(x_i^{t+1}\) as the maximizer with score \(g_{i,t}+\eta_t v_i^U(x^t)\).

We now identify the sequence with the actual declared rule by induction.
The first point belongs to \(K_a\subseteq\ri X\), where the original gradient
exists. Its directional slopes satisfy \eqref{eq:dp}. The certified slopes \eqref{eq:d} match
them on every probe ray, and the dense directions span the tangent space.
Therefore \(g_{i,1}\) is the original normalized gradient. The optimizer
just identified is the original \eqref{eq:md} optimizer and is unique, so \(x^2\)
is the actual next iterate. The canonical admissibility assumption puts
\(x^2\) in the original declared gradient domain. At this known actual point,
\eqref{eq:dp}, the certified ray identities, and dense spanning again identify
\(g_{i,2}\) with the normalized original gradient. The same argument proves
the assertion at every date.

This induction uses the original gradient-domain preservation; it does not
deduce that arbitrary points admitted by a ray certificate belong to that
domain. It permits the original finite-domain G\^ateaux interpretation as well
as the Fr\'echet interpretation, and does not require the regularizer to be
finite near an allowed boundary iterate.

Finally all marked inequalities \eqref{eq:f5} survive in the limit, so
\(R^U(x^{T_j})\ge1/b\) at infinitely many dates. This trajectory does not
approach the Nash set. We have proved the equivalence of assertions 1 and 2.

\textbf{The ordinal version.}
For a countable well-founded tree, well-founded recursion defines
\[
\rho(v)=\sup\{\rho(w)+1:w\text{ is an immediate child of }v\},
\qquad \sup\varnothing=0.
\]
Each rank is a countable ordinal, since each node has at most countably many
children and a countable supremum of countable ordinals is countable.
The tree has countably many nodes, so all its ranks are bounded by a single
countable ordinal. The ranks strictly decrease along extensions. Conversely
an infinite branch would give an infinite strictly decreasing sequence of
ordinals, which is impossible. This proves assertions 2 and 3 equivalent and
completes Theorem~\ref{main}.

\section{Why this does not change the problem's scope}

The original regularizers and every original step are retained. No entropy,
summability, small-step, strong-convexity modulus, potential structure, gradient
Lipschitz bound, or boundary-gradient continuity is added. The interior initial
set is covered by all compact cores, rather than replaced by one selected core.
The update is unique by hypothesis, so there is no additional tie-breaking
convention. All players use the same old profile in the simultaneous payoff
vectors of \eqref{eq:f4}. The target is distance to the Nash set, not action convergence,
time-averaged convergence, or convergence to one equilibrium.

The condition is a payoff-structure criterion in a closed-certificate sense:
it consists of explicitly specified finite constraints with payoff-polynomial
coefficients, followed by well-foundedness or an ordinal ranking. Real
trajectory coordinates appear only as finite feasibility variables and are
eliminated from the countable tree itself. Neither exact gradient evaluation
nor infinite orbit quantification defines node validity. The supporting
regularizer probes and ray-specific secant certificates are needed because
ordinary compactness applied to unclosed update graphs would be insufficient.

This is an infinitary representation. An arbitrary fixed regularizer or step
sequence need not have an effective finite encoding. Even decidability of
each finite node test would not decide well-foundedness of the entire tree.
Accordingly no recognition algorithm or running-time bound is claimed.
For general regularizers these finite tests need not be polynomial or
semialgebraic: fixed epigraph membership is part of the rule's data.

The tree has countably many possible waiting times and secant-denominator labels.
In general, countable well-founded trees need not have finite height; the
criterion itself imposes no uniform finite horizon. No uniform bound on the
number of bad dates or natural-number ranking is inferred. Countable ordinals
retain the full pointwise-in-initialization convergence property.

The earlier summable-entropy classification and fixed-game obstructions settle
restricted cases only. The present equivalence is not proved by extending those
restricted claims; it independently covers the canonical general fixed rule.

\section{A boundary test for pointwise ray certificates}

On the three-action simplex write the first two probabilities as \(x,y\),
so \(x,y\ge0\) and \(x+y\le1\). Consider
\begin{equation}
h(x,y)=x^2+y^2+\frac{y^2}{\sqrt x}\quad(x>0),\qquad
h(0,0)=0,\qquad h(0,y)=+\infty\quad(y>0).
\label{eq:b}
\end{equation}
It is proper and lower semicontinuous. The Hessian of \(y^2/\sqrt x\)
for \(x>0\) has nonnegative diagonal entries and determinant
\(y^2/(2x^3)\ge0\), so this term is convex there. Nonnegativity gives lower semicontinuity at the origin, and the term
\(y^2/\sqrt{x}\) tends to infinity when \(x\downarrow0\) and \(y\) approaches
a positive boundary value. For a segment from the origin to a finite-value
point with \(x>0\), \(h(tx,ty)=t^2(x^2+y^2)+t^{3/2}y^2/\sqrt{x}\le t h(x,y)\)
for \(0\le t\le1\); the interior Hessian handles all other finite-value
segments, so the stated extension is convex.
Adding \(x^2+y^2\) makes \(h\) strictly convex, and it is smooth on the
relative interior.

At the vertex \((0,0)\), every finite-domain ray has linear directional
gradient zero: for a fixed interior probe \(q=(a,b,1-a-b)\), put
\(d=q-(0,0,1)=(a,b,-a-b)\), using the ambient Euclidean norm. Then
\[
\frac{h(\theta a,\theta b)}{\theta}
=\theta(a^2+b^2)+\sqrt\theta\,\frac{b^2}{\sqrt a}
\longrightarrow0.
\]
Yet \(h(r^4,r)=r^8+r^2+1\) does not even tend to \(h(0,0)\).
A uniform Fr\'echet certificate would exclude this declared weak boundary
gradient. The repaired certificate includes it: take \(c=0\), \(g=0\).
The supporting probes are nonnegative, and, for each fixed \(q\) and \(k\),
choose a sufficiently large integer \(N\) so that
\[
N h(a/N,b/N)=\frac{a^2+b^2}{N}+\frac{b^2}{\sqrt{aN}}
\le\frac{\|d\|}{k}.
\]
This is exactly \eqref{eq:g2}, with no uniformity over probes required.

Such a vertex is not merely an irrelevant initialization. For one player,
\(\eta_1=1\), payoff vector \((7/8,-7/2,0)\), and the interior profile
\((1/16,1/4,11/16)\), the regularizer gradient representative is
\((-15/8,5/2,0)\). The next optimization score is \((-1,-1,0)\),
so its unique optimizer is \((0,0,1)\): every other point has strictly
negative score minus the nonnegative regularizer. If the declared boundary
gradient is the finite-domain linear directional gradient zero, the following
unit-step optimizer is \((7/16,0,9/16)\). The ray repair permits this exact
prefix. We do not use this test as a replacement for the full theorem or claim
a separate all-initialization convergence result from this prefix.

\section{References and roles}

\begin{thebibliography}{9}
\bibitem{potentialness}
Martin Bichler, Davide Legacci, Panayotis Mertikopoulos,
Matthias Oberlechner, and Bary Pradelski.
\emph{Characterizing the Convergence of Game Dynamics via Potentialness},
arXiv:2503.16285v1, submitted 20 March 2025.
\url{https://arxiv.org/html/2503.16285v1}.
Sections 3--4 provide game/potentialness context; Section 6 states the general
classification agenda. This reference supplies the source relationship, not
the closed-certificate characterization proved here.
\bibitem{nash}
J. F. Nash, Jr., \emph{Equilibrium points in n-person games},
Proceedings of the National Academy of Sciences USA \textbf{36}(1), 48--49
(1950), DOI: 10.1073/pnas.36.1.48.
\url{https://pmc.ncbi.nlm.nih.gov/articles/PMC1063129/}.
Used only for nonemptiness of the finite game's Nash set.
\end{thebibliography}

The interior-ray and optimizer certificates, the closed tree construction, compact
extraction, and the exact equivalence are proved above. No Lean formalization,
simulation-based proof, novelty certificate, or efficient decision procedure is
claimed.
\end{document}
